\documentclass[a4paper]{article}
\usepackage[margin=16mm]{geometry}
\usepackage{xcolor}
\usepackage{graphicx}
\graphicspath{{cover/assets/}}
\usepackage{tikz}
\usetikzlibrary{calc,shadows.blur,positioning}
\usepackage{fontspec}

% \usepackage{xeCJK}
% \setCJKmainfont{Source Han Serif SC}
% \setCJKsansfont{Source Han Sans SC}
\setmainfont{TeX Gyre Heros}
\newfontfamily\TitleFont{TeX Gyre Heros}[Scale=1.15]

\pagestyle{empty}

% 主题色
\definecolor{team}{HTML}{0EA5B5}
\definecolor{soft}{HTML}{EDF2F7}

% 变量
\newcommand{\SchoolName}{HUST}
\newcommand{\CoverTitle}{ICPC Template}
\newcommand{\Version}{v2.0}
\newcommand{\CoverDate}{\today}

% 图片文件（现在走 \graphicspath，保持这些文件名即可）
\newcommand{\SchoolLogo}{hustclogo.pdf}
\newcommand{\HeroImage}{AqoursYou.png}
\newcommand{\AvatarA}{blueMoon137.jpg}
\newcommand{\AvatarB}{sqsfx.jpg}
\newcommand{\AvatarC}{Arkweedy.jpg}

\newcommand{\HandleA}{blueMoon137}
\newcommand{\HandleB}{sqsfx}
\newcommand{\HandleC}{Arkweedy}

% —— 圆形头像样式（1 个参数=直径），关键：clip + keepaspectratio
% —— 圆形头像样式（#1 = 直径）
\tikzset{
  avatar/.style n args={1}{
    circle,
    minimum size=#1,
    inner sep=0pt,
    outer sep=0pt,
    draw=team!70!black,
    line width=0.6pt,
    path picture={
      \begin{scope}
        % 以节点圆心为中心进行圆形裁剪
        \clip (path picture bounding box.center) circle[radius=#1/2];
        % 关键：只给 height，保持等比；并且用中心锚点
        \node[anchor=center, inner sep=0pt]
          at (path picture bounding box.center)
          {\includegraphics[height=#1]{\AvatarTemp}};
      \end{scope}
    }
  }
}



\begin{document}
\begin{tikzpicture}[remember picture, overlay]

% 背景
\node[fill=soft, rounded corners=14pt, inner sep=0pt,
      minimum width=\paperwidth-10mm, minimum height=\paperheight-10mm,
      anchor=center] at (current page.center) {};

% 顶部
\node[anchor=west, xshift=8mm, yshift=-8mm, text=team, font=\Large\bfseries]
  at (current page.north west) {\SchoolName};
\node[anchor=north east, xshift=-10mm, yshift=-8mm]
  at (current page.north east) {\includegraphics[height=14mm]{\SchoolLogo}};

% 标题
\node[anchor=north, yshift=-22mm, text=black,
      font=\TitleFont\bfseries\fontsize{28}{30}\selectfont]
  at (current page.north) {\CoverTitle};

% 中央大图
\def\heroW{0.70\paperwidth}
\def\heroH{0.33\paperheight}
\node[rounded corners=16pt, draw=black!10, line width=0.6pt,
      inner sep=0pt, minimum width=\heroW, minimum height=\heroH,
      blur shadow={shadow xshift=1pt, shadow yshift=-1pt, opacity=0.08}]
      (hero) at ($(current page.center)+(0mm,10mm)$) {};
\begin{scope}
  \clip (hero.south west) rectangle (hero.north east);
  \node at (hero.center) {\includegraphics[width=\heroW,height=\heroH,keepaspectratio]{\HeroImage}};
\end{scope}

% 列坐标
% 左侧头像列的 x
\path ($(current page.west)+(0.18\paperwidth,0)$) coordinate (LeftCol);
% 右侧“网名列”的 x（靠右，右对齐）
\path ($(current page.east)+(-0.10\paperwidth,0)$) coordinate (RightCol);

% 第一行 y：放在主图下方 14mm，避免重叠
\path ($(hero.south)+(0,-14mm)$) coordinate (RowOneY);

\def\rowGap{20mm}   % 行间距
\def\avaSize{18mm}  % 头像直径
\def\nameWidth{0.30\paperwidth} % 右侧名字文本宽度

% ---------- 第一位 ----------
\def\AvatarTemp{\AvatarA}
\node[avatar=\avaSize, anchor=west] (ava1) at (LeftCol |- RowOneY) {};
\node[anchor=east, font=\large\bfseries, text width=\nameWidth, align=right]
     at (RightCol |- ava1) {\HandleA};

% ---------- 第二位 ----------
\path (ava1) ++(0,-\rowGap) coordinate (RowTwoY);
\def\AvatarTemp{\AvatarB}
\node[avatar=\avaSize, anchor=west] (ava2) at (LeftCol |- RowTwoY) {};
\node[anchor=east, font=\large\bfseries, text width=\nameWidth, align=right]
     at (RightCol |- ava2) {\HandleB};

% ---------- 第三位 ----------
\path (ava2) ++(0,-\rowGap) coordinate (RowThreeY);
\def\AvatarTemp{\AvatarC}
\node[avatar=\avaSize, anchor=west] (ava3) at (LeftCol |- RowThreeY) {};
\node[anchor=east, font=\large\bfseries, text width=\nameWidth, align=right]
     at (RightCol |- ava3) {\HandleC};

% 底部版本/日期
\node[anchor=south west, xshift=8mm, yshift=8mm, text=black!70]
  at (current page.south west) {\small Version: \textbf{\Version}};
\node[anchor=south east, xshift=-8mm, yshift=8mm, text=black!70]
  at (current page.south east) {\small Date: \textbf{\CoverDate}};

\end{tikzpicture}
\end{document}
